\section*{第 \ref{ch:g8:combustion-and-fuels} 章——燃烧：燃料、产物与温室气体}

\begin{solution}{exo:g8:combustion-and-fuels:1}
\ce{C + O2 -> CO2}.
\end{solution}

\begin{solution}{exo:g8:combustion-and-fuels:2}
氧气充足时，碳全部生成二氧化碳（完全燃烧）。
氧气不足时，部分碳生成一氧化碳或炭黑（不完全燃烧）。
\end{solution}

\begin{solution}{exo:g8:combustion-and-fuels:3}
它无色、无臭、无味，人们不易察觉，
但它有毒。
\end{solution}

\begin{solution}{exo:g8:combustion-and-fuels:4}
化石燃料：煤、天然气、汽油。可再生燃料：木材、沼气、
由甜菜制得的乙醇。
\end{solution}

\begin{solution}{exo:g8:combustion-and-fuels:5}
\ce{C3H8 + 5O2 -> 3CO2 + 4H2O}.
\end{solution}

\begin{solution}{exo:g8:combustion-and-fuels:6}
碳：2 \ce{CO2}；氢：3 \ce{H2O}；右边的氧为 $4 + 3 =
7$，其中 1 个来自乙醇，因此需 3 \ce{O2}：
\ce{C2H6O + 3O2 -> 2CO2 + 3H2O}.
\end{solution}

\begin{solution}{exo:g8:combustion-and-fuels:7}
炭黑，即碳。燃烧不完全：
火焰缺氧。
\end{solution}

\begin{solution}{exo:g8:combustion-and-fuels:8}
1980 年约为 339 ppm，2000 年约为 370 ppm，上升约 30 ppm。
\end{solution}

\begin{solution}{exo:g8:combustion-and-fuels:9}
木材中的碳是树木几年或几十年前从空气中
吸收的，燃烧时这些二氧化碳又被释放回空气中。煤中的碳
在地下封存了数百万年，燃烧时会向空气中
增加新的二氧化碳。
\end{solution}

\begin{solution}{exo:g8:combustion-and-fuels:10}
\ce{2C + O2 -> 2CO}.
\end{solution}

\begin{solution}{exo:g8:combustion-and-fuels:11}
$(427 - 316)/316 = 111/316 \approx 0.35$，约为 \qty{35}{\percent}。
\end{solution}

\begin{solution}{exo:g8:combustion-and-fuels:12}
起初氧气充足：\ce{CH4 + 2O2 -> CO2 + 2H2O}。随着
密闭房间内的氧气被消耗，火焰因空气不足而燃烧：
\ce{2CH4 + 3O2 -> 2CO + 4H2O}。一氧化碳无味且有毒，
在房间内逐渐积累。
\end{solution}

\begin{solution}{pb:g8:combustion-and-fuels:1}
\textbf{1.} \ce{C3H8 + 5O2 -> 3CO2 + 4H2O}.

\textbf{2.} 左边：3 个 C，8 个 H，10 个 O。右边：3 个 C，$4 \times 2 = 8$ 个 H，
$3 \times 2 + 4 = 10$ 个 O。配平。

\textbf{3.} \ce{2C3H8 + 7O2 -> 6CO + 8H2O}.

\textbf{4.} 完全燃烧时，每个丙烷分子需 5 个氧分子，
不完全燃烧时则为 $7/2 = 3.5$ 个。

\textbf{5.} 密闭帐篷内的空气很快缺氧，
火焰得不到足够的氧气，燃烧不完全，火焰呈黄色。

\textbf{6.} 一氧化碳无色、无臭（也无味）。

\textbf{7.} GHS02 为火焰标志：极易燃。GHS04 为
气瓶标志：气体受压存放。

\textbf{8.} 例如：只能在户外使用，切勿在帐篷或
封闭房间中使用；远离一切可燃物；更换
气罐时须远离任何明火。

\textbf{9.} 重 $132 \div 44 = 3$ 倍。

\textbf{10.} 来自空气中的氧，它最终进入二氧化碳
（以及水）中。

\textbf{11.} 属于化石燃料（来自天然气或石油），燃烧
会向空气中增加二氧化碳。

\textbf{12.} 共产生二氧化碳
$450 \times 3 = \qty{1350}{g} = \qty{1.35}{kg}$。
\end{solution}
